\documentclass[11pt,reqno]{amsart}

\usepackage[T1]{fontenc}
\usepackage{lmodern}
\usepackage{amsmath,amssymb,mathtools}
\usepackage{booktabs,array}
\usepackage{enumitem}
\usepackage{microtype}
\usepackage{needspace}
\usepackage[colorlinks=true,linkcolor=blue,citecolor=blue,urlcolor=blue]{hyperref}

\newtheorem{theorem}{Theorem}[section]
\newtheorem{proposition}[theorem]{Proposition}
\newtheorem{lemma}[theorem]{Lemma}
\newtheorem{corollary}[theorem]{Corollary}
\theoremstyle{definition}
\newtheorem{definition}[theorem]{Definition}
\newtheorem{example}[theorem]{Example}
\theoremstyle{remark}
\newtheorem{remark}[theorem]{Remark}

\DeclareMathOperator{\Aut}{Aut}
\DeclareMathOperator{\supp}{supp}
\DeclareMathOperator{\tr}{tr}
\newcommand{\C}{\mathbf C}
\newcommand{\F}{\mathbf F}
\newcommand{\Q}{\mathbf Q}
\newcommand{\T}{\mathbf T}
\newcommand{\Z}{\mathbf Z}

\title[Floquet sparse obstruction]{Floquet Sparse Obstruction:
Exact Exclusion Through Support Six in the Period-Three Hamming Fibre}
\author{KEIO}
\date{September 4, 2026}

\hypersetup{
  pdftitle={Floquet Sparse Obstruction: Exact Exclusion Through Support Six in the Period-Three Hamming Fibre},
  pdfauthor={KEIO},
  pdfsubject={Floquet isospectrality and sparse diagonal perturbations},
  pdfkeywords={Floquet isospectrality, Hamming graph, sparse potential, Groebner basis, exact computation}
}

\subjclass[2020]{Primary 47B36; Secondary 05C50, 13P10, 35P05}
\keywords{Floquet isospectrality, discrete Schr\"odinger operator,
Hamming graph, sparse diagonal perturbation, exact computation,
Gr\"obner basis}

\begin{document}

\begin{abstract}
Let $\Delta$ be the adjacency-form discrete Laplacian on $\Z^d$, and let
$V\colon\Z^d\to\C$ be $(3\Z)^d$-periodic.  We study complex potentials for
which $\Delta+V$ and $\Delta$ are Floquet isospectral.  At zero
quasimomentum the Floquet matrix is the adjacency matrix of the Hamming graph
$H(d,3)$.  Trace identities at this single fibre show, by a completely
human-checkable argument, that every nonzero isospectral diagonal shift of
this fibre has at least five nonzero entries for every $d\geq1$.  For
$d\geq3$, equality of the first five traces on a support of size at most four
has a unique nonzero form: an induced claw with centre value $-1$ and leaf
values $1,i,-i$.  Its sixth trace has the nonzero defect $-12(2d-5)$, so the
sixth moment is sharp for this sparse trace test.

For the smallest unresolved rectangular period, $(3\Z)^3$, we strengthen the
lower bound to seven by an exact computer-assisted proof.  The resolvent of
$H(3,3)$ depends only on Hamming distance.  The matrix determinant lemma
therefore reduces supports of sizes five and six to, respectively, 104 and
320 distance-matrix types.  After imposing exact support by Rabinowitsch
extension, every resulting extension ideal has reduced Gr\"obner basis
$\{1\}$ over $\Q$.  The computation uses exact rational polynomial arithmetic; an
independent orbit and Burnside audit verifies the completeness of the type
lists, and a deterministic transcript digest records every type.  Thus every
nonzero shift of support at most six is detected already at the periodic
fibre.  This closes the support-at-most-six stratum of the open $(3,3,3)$
problem, but does not decide whether a nonzero potential of unrestricted
support exists.
\end{abstract}

\maketitle

\section{Introduction}

For $n\in\Z^d$ put
\[
 (\Delta f)(n)=\sum_{\lVert n-m\rVert_1=1} f(m).
\]
Let $V$ be $(3\Z)^d$-periodic and identify a fundamental cell with
$\mathcal F_d=\{0,1,2\}^d$.  Put
$\T=\{z\in\C:|z|=1\}$.  For $z=(z_1,\ldots,z_d)\in\T^d$, the Floquet
matrix $L_V(z)$ is defined by
\[
 (L_V(z)f)(n)=V(n)f(n)+
 \sum_{j=1}^d\bigl(\eta_j^+(n)f(n+e_j)
              +\eta_j^-(n)f(n-e_j)\bigr),
\]
where the arguments of $f$ are reduced modulo $3$ and
\[
 \eta_j^+(n)=\begin{cases}z_j,&n_j=2,\\1,&n_j\neq2,\end{cases}
 \qquad
 \eta_j^-(n)=\begin{cases}z_j^{-1},&n_j=0,\\1,&n_j\neq0.\end{cases}
\]
We say that $V$ is \emph{Floquet isospectral to zero} when
\begin{equation}\label{eq:floquet-definition}
 \det(xI-L_V(z))=\det(xI-L_0(z))
 \qquad(z\in\T^d).
\end{equation}
The trace and spectral-invariant approach to discrete periodic inverse
problems goes back at least to Kappeler~\cite{kappeler1989}; see also
Liu~\cite{liu2023}, Saburova's complete invariant system for periodic
graphs~\cite{saburova2026}, and the algebraic framework surveyed by Shipman
and Sottile~\cite{shipman2026}.  Real-valued potentials isospectral to zero
are rigid, whereas nonzero complex examples occur for several larger
periods.  Faust, Liu, Matos, Plute, Robinson, Tao, Tran and Zhuang gave
explicit complex examples, including a support-four construction with
values $1+i,1-i,-1+i,-1-i$ for even periods at least four~\cite{faust2024}.

Cobb, Faust and Kretschmer proved existence when at least one rectangular
period exceeds $3$, and rigidity when all periods are at most $3$ and at
most two of them equal $3$~\cite{cobb2026}.  Thus $(3,3,3)$ is the smallest
rectangular period not covered by their classification; they also record
that generating its full $27$-variable spectral-invariant ideal was beyond
the direct methods they attempted.  The general period-classification
problem and the all-period-three barrier are stated explicitly in
\cite[Problem~1]{faustideal2026}.

This paper does not solve the unrestricted $(3,3,3)$ case.  It isolates a
sparse stratum and replaces the full ideal by exact low-rank determinant
identities at one necessary Floquet fibre.  The result is an exact exclusion
through support six, together with a finite decision procedure for every
fixed Hamming-fibre support size.  Here ``exact'' refers to exhaustive type
enumeration and rational polynomial arithmetic; it does not assert that a
support-seven solution exists.

\begin{theorem}[Human Hamming-fibre obstruction]\label{thm:human}
Let $d\geq1$, let $A_d=A(H(d,3))$, and let $D$ be a complex diagonal
matrix.  If
$A_d+D$ and $A_d$ have the same characteristic polynomial, then
\[
 D\neq0\quad\Longrightarrow\quad
 |\supp D|\geq5.
\]
\end{theorem}

\begin{theorem}[Exact $H(3,3)$ obstruction]\label{thm:computer}
Let $A=A(H(3,3))$ and let $D$ be a complex diagonal matrix.  If $A+D$ and
$A$ have the same characteristic polynomial, then
\[
 D\neq0\quad\Longrightarrow\quad
 \boxed{|\supp D|\geq7}.
\]
The exclusion of supports five and six is computer-assisted and uses exact
arithmetic over $\Q$.
\end{theorem}

\begin{corollary}[Floquet consequence]\label{cor:floquet}
If a $(3\Z)^3$-periodic complex potential $V$ is Floquet isospectral to
zero, then $V=0$ or its support in a fundamental cell has size at least
seven.
\end{corollary}

\begin{corollary}[Non-vacuity at the periodic fibre]\label{cor:nonvacuity}
There exists a nonzero complex diagonal matrix $D$ such that $A(H(3,3))+D$
and $A(H(3,3))$ have the same characteristic polynomial, and every such
$D$ has support at least seven.
\end{corollary}

\begin{proof}
The graph $H(3,3)$ is neither edgeless nor complete.  The fixed-matrix
existence theorem of Cobb--Faust--Kretschmer
\cite[Corollary~B]{cobb2026} therefore supplies a nonzero isospectral
diagonal shift.  Its support bound is Theorem~\ref{thm:computer}.  This
fixed-fibre existence statement does not imply simultaneous isospectrality
at all Floquet phases.
\end{proof}

The proof of Theorem~\ref{thm:human} uses only traces through order six and a
nine-row graph table.  Theorem~\ref{thm:computer} uses the determinant lemma,
the Hamming association scheme and Rabinowitsch-extended coefficient ideals.  The
calculation is finite because the zero-quasimomentum resolvent is a function
of pairwise Hamming distance.  No floating-point arithmetic, numerical root
test or tolerance enters the calculation.

We are not aware of a previously stated support lower bound of this form.
That literature assessment is not a priority certificate; the claim should
be revisited after expert feedback.

\section{The zero Floquet fibre}

Write
\[
 \mathcal F_d=(\Z/3\Z)^d.
\]
At $z=(1,\ldots,1)$ the free Floquet matrix is
\[
 A_d=A(H(d,3)),
 \qquad H(d,3)=K_3\mathbin{\square}\cdots\mathbin{\square}K_3.
\]
For a diagonal matrix $D$ indexed by $\mathcal F_d$, set
\[
 \supp D=\{a\in\mathcal F_d:D_{aa}\neq0\}.
\]
Let $D$ be any complex diagonal matrix such that $A_d+D$ and $A_d$ have
the same characteristic polynomial.  Newton identities then give
\begin{equation}\label{eq:trace-equality}
 \tr(A_d+D)^k=\tr A_d^k\qquad(k\geq1).
\end{equation}
For a Floquet-isospectral potential one takes
$D=\operatorname{diag}(V|_{\mathcal F_d})$; equation
\eqref{eq:floquet-definition} implies the required finite-matrix hypothesis.
Only this necessary zero-quasimomentum condition will be used.
These power traces are the single-fibre specialization of the spectral
invariants studied more generally in~\cite{saburova2026}.

The following elementary local identities are immediate from the product
graph:
\begin{equation}\label{eq:local-walks}
 (A_d)_{xx}=0,\qquad
 (A_d^2)_{xx}=2d,\qquad
 (A_d^3)_{xx}=2d,
\end{equation}
and
\begin{equation}\label{eq:common-neighbour}
 (A_d^2)_{xy}=1\qquad(x\sim y).
\end{equation}
Every edge lies in a unique triangle: its endpoints differ in one coordinate,
and the third symbol in that coordinate gives the unique common neighbour
adjacent to both within the same coordinate fibre.

\begin{lemma}[The first three moments]\label{lem:first-three}
Let $v_1,\ldots,v_r$ be the nonzero entries of $D$, and put
$p_j=\sum_{a=1}^r v_a^j$.  If \eqref{eq:trace-equality} holds for
$k=1,2,3$, then
\[
 p_1=p_2=p_3=0.
\]
\end{lemma}

\begin{proof}
The first trace gives $p_1=0$.  Since $A_d$ has zero diagonal,
\[
 \tr(A_d+D)^2-\tr A_d^2=p_2.
\]
For the third power, cyclicity of trace and \eqref{eq:local-walks} give
\[
 \tr(A_d+D)^3-\tr A_d^3
 =p_3+3\tr(A_d^2D)=p_3+6dp_1=p_3.
\]
\end{proof}

\begin{corollary}[Value-level obstruction]\label{cor:value-levels}
Every nonzero isospectral diagonal shift of $A_d$ assumes at least four
distinct nonzero values on its support.
\end{corollary}

\begin{proof}
Suppose the distinct nonzero values are $a_1,\ldots,a_s$, with positive
multiplicities $m_1,\ldots,m_s$, where $s\leq3$.  Lemma~\ref{lem:first-three}
gives
\[
 \sum_{j=1}^s m_j a_j^k=0\qquad(k=1,2,3).
\]
The first $s$ equations have coefficient matrix
$(a_j^k)_{1\leq k,j\leq s}$, whose determinant is a nonzero Vandermonde
factor because the $a_j$ are nonzero and distinct.  Hence every $m_j=0$, a
contradiction.
\end{proof}

\section{A human proof through support four}

\begin{lemma}[Newton reduction]\label{lem:newton}
No nonzero isospectral diagonal shift of $A_d$ has support at most three.
If its support has size four, then for some $\alpha\in\C^\times$ its four
values are
\begin{equation}\label{eq:four-values}
 \alpha\{1,-1,i,-i\}.
\end{equation}
\end{lemma}

\begin{proof}
Newton's identities and Lemma~\ref{lem:first-three} give
$e_1=e_2=e_3=0$ for the elementary symmetric functions of the nonzero
values.  If $r\leq3$, the product $e_r$ is therefore zero, contrary to the
definition of the support.  If $r=4$, the monic polynomial with roots
$v_1,\ldots,v_4$ has the form $X^4+e_4$, which yields
\eqref{eq:four-values} after choosing a fourth root.
\end{proof}

Let $S$ be a four-point support and write
$v_x=\alpha\omega_x$, where the labels on $S$ are the four fourth roots of
unity.  From the graph induced on $S$ define
\begin{equation}\label{eq:ET}
 E=\sum_{\{x,y\}\in E(S)}\omega_x\omega_y,
 \qquad
 T=\sum_{\{x,y\}\in E(S)}
 \omega_x\omega_y(\omega_x+\omega_y).
\end{equation}

\begin{lemma}[Trace four and trace five]\label{lem:trace45}
Assume that $D$ has four-point support and write its values as
$v_x=\alpha\omega_x$ as above.  Under the first-three-moment relations,
\begin{align}
 \tr(A_d+D)^4-\tr A_d^4
   &=4\alpha^2(E+\alpha^2),\label{eq:trace4}\\
 \tr(A_d+D)^5-\tr A_d^5
   &=5\alpha^2(2E+\alpha T).\label{eq:trace5}
\end{align}
Consequently a four-point candidate must satisfy
\begin{equation}\label{eq:ET-condition}
 E\neq0,\qquad T^2+4E=0,
 \qquad \alpha=-\frac{2E}{T}.
\end{equation}
\end{lemma}

\begin{proof}
The graph is vertex-transitive, so $(A_d^j)_{xx}$ is independent of $x$;
hence every cyclic word containing exactly one $D$ is a multiple of $p_1$.
 At order four, cyclic word classes give
\[
 \tr(A_d+D)^4-\tr A_d^4
 =\tr D^4+2\tr(A_dD A_dD),
\]
after the terms containing $p_1,p_2,p_3$ are removed.  Since
$\tr D^4=4\alpha^4$ and
$\tr(A_dD A_dD)=2\alpha^2E$, this is \eqref{eq:trace4}.  At order five, the only surviving cyclic
classes are $A_d^2DA_dD$ and $A_dDA_dD^2$.  Using
$(A_d^2)_{xy}=1$ on an edge, their total contribution is
\[
 10\alpha^2E+5\alpha^3T,
\]
while $\tr D^5=\alpha^5\sum\omega_x=0$.  This is
\eqref{eq:trace5}.  If both differences vanish, then
$\alpha^2=-E$, so $E\neq0$, and $\alpha T=-2E$.  Eliminating $\alpha$
gives \eqref{eq:ET-condition}.
\end{proof}

There are eleven unlabelled simple graphs on four vertices.  The unique
triangle property excludes $K_4$ and the diamond.  Fixing one root label to
$1$ leaves six permutations of the other labels.  This normalization is
lossless: for any fourth root $\zeta$, replacing
$(\alpha,\omega_x)$ by $(\alpha/\zeta,\zeta\omega_x)$ preserves every
value $v_x$ and the trace conditions.  Direct substitution in
\eqref{eq:ET-condition} gives the complete table below.

\begin{table}[htbp]
\centering
\caption{The complete trace-four/five test after normalizing one root label
to $1$.  Graphs $K_4$ and the diamond cannot occur in $H(d,3)$.}
\label{tab:four-point}
\begin{tabular}{lrr}
\toprule
Induced graph & labelings with $E\neq0$ & survive traces $4,5$\\
\midrule
$4K_1$       &0&0\\
$K_2+2K_1$   &6&0\\
$P_3+K_1$    &4&0\\
$K_3+K_1$    &6&0\\
$K_{1,3}$    &6&6\\
$2K_2$       &4&0\\
$P_4$        &6&0\\
paw          &4&0\\
$C_4$        &4&0\\
\bottomrule
\end{tabular}
\end{table}

For the claw, let $c$ be the root label at the centre.  The leaf labels are
the other three roots, whence
\[
 E=-c^2,\qquad T=-2c^3.
\]
The trace equations force $\alpha c=-1$.  Thus the centre has value $-1$
and the leaf values are $\{1,i,-i\}$.  Independent leaves adjacent to a
common vertex must lie in three distinct coordinate directions.  Such a claw
exists exactly when $d\geq3$, and all such configurations are conjugate
under $\Aut H(d,3)$.

\begin{lemma}[The sixth-trace defect]\label{lem:trace6}
For the formal claw candidate just described,
\begin{equation}\label{eq:sixth-defect}
 \tr(A_d+D)^6-\tr A_d^6=-12(2d-5).
\end{equation}
\end{lemma}

\begin{proof}
Group the words of length six by cyclic rotation.  After the all-$A$ word is
removed, the exact contributions are as follows.  They can be checked from
the distance-walk recurrence
\[
 w_{\ell+1}(h)
 =h\,w_\ell(h-1)+h\,w_\ell(h)+2(d-h)w_\ell(h+1),
 \qquad w_0(h)=\mathbf 1_{h=0},
\]
where $w_\ell(h)=(A_d^\ell)_{xy}$ whenever the Hamming distance from $x$ to
$y$ is $h$, together with multiplication on the four supported vertices.
\begin{table}[htbp]
\centering
\caption{Cyclic word contributions for the formal claw.}
\label{tab:trace-six}
\begin{tabular}{lrr@{\qquad}lrr}
\toprule
word & mult. & trace & word & mult. & trace\\
\midrule
$A^5D$     &6&0          &$A^4D^2$   &6&0\\
$A^3DAD$   &6&$-12d+6$   &$A^2DA^2D$ &3&6\\
$A^3D^3$   &6&0          &$A^2DAD^2$ &6&2\\
$A^2D^2AD$ &6&2          &$ADADAD$    &2&0\\
$A^2D^4$   &6&$8d$       &$ADAD^3$    &6&$-2$\\
$AD^2AD^2$ &3&$-2$       &$AD^5$      &6&0\\
$D^6$      &1&0          &&&\\
\bottomrule
\end{tabular}
\end{table}
The entries follow by multiplying on the four supported vertices, using
\eqref{eq:local-walks}--\eqref{eq:common-neighbour} and the fact that the
three leaves occupy distinct coordinate directions.  Multiplicity times
trace, summed over Table~\ref{tab:trace-six}, is $-12(2d-5)$.  The ancillary
verifier performs this calculation symbolically in $d$ from the displayed
distance recurrence and asserts all thirteen rows separately.
\end{proof}

\begin{proposition}[Sharp five-moment classification]\label{prop:five-moment}
Let $d\geq3$ and $0<|\supp D|\leq4$.  Then
\[
 \tr(A_d+D)^k=\tr A_d^k\qquad(1\leq k\leq5)
\]
if and only if the support is an induced claw and, after permuting its
leaves, the values are
\[
 D_{\mathrm{centre}}=-1,
 \qquad
 (D_{\mathrm{leaf},1},D_{\mathrm{leaf},2},D_{\mathrm{leaf},3})=(1,i,-i).
\]
Every such shift has sixth-trace defect $-12(2d-5)\neq0$.  In particular,
order six is sharp for detecting all nonzero shifts of support at most four
by consecutive power traces.
\end{proposition}

\begin{proof}
Necessity is Lemma~\ref{lem:newton}, Table~\ref{tab:four-point}, and the
calculation following that table.  Conversely, the displayed claw has
$p_1=p_2=p_3=0$, $E=-c^2$, $T=-2c^3$, and $\alpha c=-1$, so
Lemma~\ref{lem:trace45} verifies the first five traces.  The last assertion
is Lemma~\ref{lem:trace6}.
\end{proof}

\begin{proof}[Proof of Theorem~\ref{thm:human}]
Lemma~\ref{lem:newton} excludes support at most three in every dimension.
For $d=1$ this exhausts the three-vertex cell.  For $d=2$, the neighbours of
a vertex induce two disjoint copies of $K_2$, one for each coordinate
direction, so they contain no independent three-set and hence no induced
claw.  Table~\ref{tab:four-point} therefore excludes support four.  For
$d\geq3$, the same table leaves only the formal claw and
Lemma~\ref{lem:trace6} excludes it because $-12(2d-5)\neq0$.
\end{proof}

\section{Distance resolvents of period-three Hamming graphs}

For $0\leq j,h\leq d$, define the ternary Krawtchouk polynomial
\begin{equation}\label{eq:krawtchouk}
 K_j^{(d)}(h)=\sum_{\ell=0}^j(-1)^\ell2^{j-\ell}
 \binom{h}{\ell}\binom{d-h}{j-\ell},
\end{equation}
with the usual convention for binomial coefficients outside their range,
and put $\theta_j=2d-3j$.

\begin{proposition}[Hamming resolvent]\label{prop:general-resolvent}
For vertices $a,b\in\mathcal F_d$ at Hamming distance $h$,
\begin{equation}\label{eq:general-resolvent}
 ((xI-A_d)^{-1})_{ab}
 =3^{-d}\sum_{j=0}^d\frac{K_j^{(d)}(h)}{x-\theta_j}.
\end{equation}
In particular, after multiplication by the minimal polynomial
\[
 m_d(x)=\prod_{j=0}^d(x-\theta_j),
\]
the matrix
\[
 N_d(x)=m_d(x)(xI-A_d)^{-1}
\]
has polynomial entries over $\Q$, each determined only by $h$.
\end{proposition}

\begin{proof}
The characters of $\mathcal F_d$ diagonalize $A_d$.  A character with
exactly $j$ nontrivial coordinates has eigenvalue $\theta_j$; summing the
rank-one character projectors of that weight gives kernel
$3^{-d}K_j^{(d)}(h)$.  The spectral decomposition yields
\eqref{eq:general-resolvent}.  The eigenvalues $\theta_0,\ldots,\theta_d$
are distinct, so their product is the minimal polynomial and clears every
denominator.
\end{proof}

\subsection{The case \texorpdfstring{$H(3,3)$}{H(3,3)}}

From now on $d=3$ and $A=A_3$.  Its distinct eigenvalues are
$6,3,0,-3$, so its minimal polynomial is
\begin{equation}\label{eq:m}
 m(x)=x(x-6)(x-3)(x+3).
\end{equation}
For vertices $a,b\in\F_3^3$, let $\delta(a,b)$ be their Hamming distance and
define
\begin{align}
 n_0(x)&=x^3-6x^2-3x+24,\nonumber\\
 n_1(x)&=x(x-5),\nonumber\\
 n_2(x)&=2(x-3),\label{eq:numerators}\\
 n_3(x)&=6.\nonumber
\end{align}
Let $N(x)$ be the $27$-by-$27$ matrix with
$N_{ab}(x)=n_{\delta(a,b)}(x)$.

\begin{proposition}[Distance resolvent]\label{prop:resolvent}
One has the polynomial matrix identity
\begin{equation}\label{eq:resolvent-identity}
 (xI-A)N(x)=m(x)I.
\end{equation}
Consequently
\[
 (xI-A)^{-1}_{ab}=\frac{n_{\delta(a,b)}(x)}{m(x)}
\]
away from the spectrum of $A$.
\end{proposition}

\begin{proof}
Both sides lie in the Bose--Mesner algebra of the Hamming association
scheme.  Multiplication by $A$ counts one-step changes from each distance
class.  Substitution of the four polynomials in \eqref{eq:numerators} into
these four distance recurrences gives \eqref{eq:resolvent-identity}.
The ancillary verifier also checks all $27^2$ entries symbolically over
$\Z[x]$ before any support calculation begins.
\end{proof}

\section{A finite Rabinowitsch ideal for a prescribed support}

Fix $S\subset\F_3^3$ of size $r$.  Let $U_S$ be the $27$-by-$r$ coordinate
embedding, write $D_S=\operatorname{diag}(v_0,\ldots,v_{r-1})$, set
$D=U_SD_SU_S^{\mathsf T}$, and let $N_S(x)=U_S^{\mathsf T}N(x)U_S$.

\begin{proposition}[Determinant-lemma reduction]\label{prop:detlemma}
The zero-fibre matrices $A+D$ and $A$ have the same characteristic
polynomial if and only if the following polynomial vanishes identically in
$x$:
\begin{equation}\label{eq:spectral-polynomial}
 F_S^{(r)}(x;v)=
 \sum_{\varnothing\neq T\subseteq\{0,\ldots,r-1\}}
 (-1)^{|T|}m(x)^{r-|T|}
 \left(\prod_{j\in T}v_j\right)
 \det N_S(x)[T,T].
\end{equation}
\end{proposition}

\begin{proof}
For generic $x$, the matrix determinant lemma and
Proposition~\ref{prop:resolvent} give
\[
 \frac{\det(xI-A-D)}{\det(xI-A)}
 =\det\left(I_r-D_S\frac{N_S(x)}{m(x)}\right).
\]
Expand the determinant by principal minors, subtract $1$, and multiply by
$m(x)^r$.  This is \eqref{eq:spectral-polynomial}.  Since both sides are
rational functions, the resulting identity is also sufficient.
\end{proof}

Let $I_S^{(r)}$ be the ideal over $\Q$ generated by the coefficients in $x$
of $F_S^{(r)}$ after the trace equation
\begin{equation}\label{eq:trace-eliminate}
 v_{r-1}=-(v_0+\cdots+v_{r-2})
\end{equation}
has been substituted.  Exact support is imposed by a new variable $u$ and
the Rabinowitsch equation
\begin{equation}\label{eq:rabinowitsch}
 u\prod_{j=0}^{r-1}v_j-1=0,
\end{equation}
again after \eqref{eq:trace-eliminate}.  Denote the resulting extension ideal
by
\begin{equation}\label{eq:extension-ideal}
 J_S^{(r)}=
 I_S^{(r)}+\left(u\prod_{j=0}^{r-1}v_j-1\right)
 \subset\Q[u,v_0,\ldots,v_{r-2}].
\end{equation}

\begin{corollary}[Exact unit-ideal criterion]\label{cor:unit}
There is no complex zero-fibre isospectral diagonal shift with exact support
$S$ if and only if $J_S^{(r)}=(1)$.  Equivalently, any reduced
Gr\"obner basis of $J_S^{(r)}$ is $\{1\}$.
\end{corollary}

\begin{proof}
Equation \eqref{eq:rabinowitsch} is solvable precisely when every $v_j$ is
nonzero.  Proposition~\ref{prop:detlemma} identifies the remaining
coefficient equations with isospectrality.  The assertion is therefore the
weak Nullstellensatz and the standard unit-ideal criterion for a Gr\"obner
basis~\cite{cox}.
\end{proof}

The matrix $N_S$ depends, up to simultaneous row and column permutation,
only on the abstract pairwise Hamming-distance matrix of $S$.  Thus one
unit-ideal calculation per distance type is exhaustive.  This reduction is
valid specifically at zero quasimomentum; it must not be transferred to a
nonzero Floquet phase without a new argument.

\begin{proposition}[Finite exact decidability]\label{prop:finite-decision}
For every fixed pair $(d,r)$ with $1\leq r\leq3^d$, the existence of an
isospectral diagonal shift of $A(H(d,3))$ with exact support $r$ is decidable
by a finite exact computation over $\Q$, with one Rabinowitsch-extended
coefficient ideal for each realizable $r$-point Hamming-distance type.
\end{proposition}

\begin{proof}
There are finitely many $r$-subsets of $\mathcal F_d$ and hence finitely many
distance types.  Proposition~\ref{prop:general-resolvent} supplies a common
rational resolvent denominator and a numerator kernel determined by the
distance matrix.  The determinant-lemma construction above therefore gives
one ideal $J_S^{(r)}$ over $\Q$ per type.  Exact-support existence for that
type is equivalent to this ideal being proper by
Corollary~\ref{cor:unit}; a terminating Gr\"obner-basis computation decides
the question.  This is a decidability statement, not a claim of practical
complexity uniformly in $(d,r)$.
\end{proof}

\Needspace{8\baselineskip}
\section{Exact exclusion of support five}\label{sec:support-five}

\begin{lemma}[Five-point completeness]\label{lem:five-completeness}
The $\binom{27}{5}=80{,}730$ five-subsets form $105$ orbits under
\[
 \Aut H(3,3)\cong(S_3)^3\rtimes S_3,
 \qquad |\Aut H(3,3)|=1296,
\]
and these orbits determine exactly $104$ abstract pairwise Hamming-distance
types.
\end{lemma}

\begin{proof}
The primary verifier enumerates all five-subsets and minimizes the ten-entry
upper-triangular distance word over all $5!$ point relabelings.  Independently,
the completeness auditor partitions all subsets by union-find under eight
generators of $(S_3)^3\rtimes S_3$, obtains $105$ orbits, and only then
canonicalizes their distance matrices.  Both routes give $104$ types and the
same serialized-key digest displayed below.  The passage from $105$ orbits
to $104$ distance types is valid because $N_S(x)$ is determined entirely by
the pairwise Hamming-distance matrix.
\end{proof}

For every type the verifier constructs \eqref{eq:spectral-polynomial}, makes
the exact substitution $v_4=-(v_0+\cdots+v_3)$, appends
\eqref{eq:rabinowitsch}, and computes a Gr\"obner basis over $\Q$ in the
variables $u,v_0,\ldots,v_3$ using graded reverse lexicographic order.

\begin{proposition}[Support-five exclusion]\label{prop:support5}
All $104$ Rabinowitsch extension ideals have Gr\"obner basis $\{1\}$.  Hence
there is no exact-support-five zero-fibre isospectral shift on $H(3,3)$.
\end{proposition}

\begin{proof}[Exact computer-assisted proof]
Lemma~\ref{lem:five-completeness} makes the type list exhaustive.  For each
type, the verifier constructs $F_S^{(5)}$, extracts every coefficient in
$x$, performs the trace substitution, and constructs $J_S^{(5)}$ from
\eqref{eq:extension-ideal}.  Exact arithmetic over $\Q$ returns the reduced
basis $\{1\}$ for each of the $104$ ideals.  Corollary~\ref{cor:unit} excludes
exact support five for every type.
\end{proof}

The sorted distance keys, serialized as comma-separated entries with one key
per line, have SHA-256 digest
\begin{center}
{\small\texttt{222cdf7f5c212970fb39bf5e723bfcf29dfdfa72c6bd77ab9f25c10e5686cea9}}.
\end{center}
The reference run used Python 3.12.13 and SymPy 1.14.0 and returned
$104/104$ unit ideals.

\section{Exact exclusion of support six}\label{sec:support-six}

\begin{lemma}[Six-point completeness]\label{lem:six-completeness}
The six-point support counts are those in Table~\ref{tab:six-counts}.
\begin{table}[htbp]
\centering
\caption{Exact combinatorial counts for six-point supports.}
\label{tab:six-counts}
\begin{tabular}{lr}
\toprule
quantity & exact value\\
\midrule
six-subsets & $\binom{27}{6}=296{,}010$\\
automorphism group order & $1{,}296$\\
automorphism orbits & $321$\\
pairwise-distance types & $320$\\
\bottomrule
\end{tabular}
\end{table}
Exactly two nonconjugate automorphism orbits share one abstract distance
matrix; all other distance types contain one orbit.
\end{lemma}

\begin{proof}
One exact computation partitions all $296{,}010$ subset masks by union-find
under eight generators of the automorphism group.  A separate Burnside sum
over all $1{,}296$ group elements gives
\[
 \frac{1}{1296}\sum_{g\in\Aut H(3,3)}
 \#\operatorname{Fix}_{\binom{\mathcal F_3^3}{6}}(g)
 =\frac{416{,}016}{1{,}296}=321,
\]
independently confirming completeness.  Canonicalizing the fifteen-entry
distance words of the $321$ orbit representatives gives $320$ types, with
one distance word occurring twice.  The two representatives and their
distinct coordinate-fibre profiles are recorded by the ancillary auditor.
\end{proof}

The distinction in the final two rows is essential.  Exactly two
nonconjugate automorphism orbits share one abstract distance matrix.  They
may be coalesced here only because the resolvent kernel in
Proposition~\ref{prop:resolvent} is distance-dependent.

For each of the $320$ types, the support-six verifier repeats the construction
of Section~\ref{sec:support-five} with $v_5=-(v_0+\cdots+v_4)$.

\begin{proposition}[Support-six exclusion]\label{prop:support6}
All $320$ Rabinowitsch extension ideals have Gr\"obner basis $\{1\}$ over
$\Q$.  Hence no zero-fibre isospectral diagonal shift on $H(3,3)$ has
exact support six.
\end{proposition}

\begin{proof}[Exact computer-assisted proof]
Lemma~\ref{lem:six-completeness} makes the distance-type list exhaustive.
For every type, the support-six verifier constructs $F_S^{(6)}$, extracts
all coefficients in $x$, performs the trace substitution, and constructs
$J_S^{(6)}$.  All $320$ exact reduced bases over $\Q$ are $\{1\}$.
Corollary~\ref{cor:unit} therefore excludes exact support six.
\end{proof}

The exact JSONL computation transcript records one header, $320$ sorted type
records and one summary record.  Its principal digests are
\begin{table}[htbp]
\centering
\caption{Digests for the support-six release.  The first two are
deterministic mathematical digests; the last is an archival file-integrity
digest.}
\label{tab:six-digests}
\begin{tabular}{l>{\raggedright\arraybackslash}p{0.58\textwidth}}
\toprule
object & SHA-256\\
\midrule
sorted distance keys &
\shortstack[l]{\scriptsize\texttt{b7a0ff2f436adb2f9eed342cf33324b4}\\
\scriptsize\texttt{45cb118672017d5e8eb6beb619037ecb}}\\
canonical transcript &
\shortstack[l]{\scriptsize\texttt{192dddaf12e4622821aec0f58ff4b95e}\\
\scriptsize\texttt{d46e853f14333c44a665c3fb0d0ada1e}}\\
complete JSONL file &
\shortstack[l]{\scriptsize\texttt{c0c64dc9f4ca7bfd8b63cbbcd08ea47}\\
\scriptsize\texttt{368948aeb31917067cf1397d9f86c7257}}\\
\bottomrule
\end{tabular}
\end{table}
The reference run returned $320/320$ unit ideals and no nonunit index.
Types $1$, $90$, $160$ and $320$ were recomputed in fresh processes and
reproduced their stored ideal hashes.  A separate union-find/Burnside program
reproduced the $321$ orbits, the $320$ types and the distance-key digest.
The complete-file digest is expected to change on a fresh run because the
log includes process identifiers, platform data and timings; it authenticates
the distributed file only.  The canonical transcript digest omits those
ephemeral fields and is the reproducible comparison target.

\begin{proof}[Proof of Theorem~\ref{thm:computer}]
Theorem~\ref{thm:human} excludes support at most four.  Propositions
\ref{prop:support5} and~\ref{prop:support6}, through
Corollary~\ref{cor:unit}, exclude exact supports five and six.
\end{proof}

\begin{proof}[Proof of Corollary~\ref{cor:floquet}]
Full Floquet isospectrality implies equality of the characteristic
polynomials at zero quasimomentum, so Theorem~\ref{thm:computer} applies.
\end{proof}

\section{Product transfer and cylinder potentials}

Let $1\leq k<d$ and let $V_0$ be a $(3\Z)^k$-periodic potential.  Its
\emph{cylinder lift} to $\Z^d$ is
\[
 V(n_1,\ldots,n_d)=V_0(n_1,\ldots,n_k).
\]

\begin{proposition}[Cylinder equivalence]\label{prop:cylinder}
The cylinder lift $V$ is Floquet isospectral to zero on $\Z^d$ if and only
if $V_0$ is Floquet isospectral to zero on $\Z^k$.
\end{proposition}

\begin{proof}
Write $z=(z',z'')$.  In the product basis the Floquet matrix is the
Kronecker sum
\[
 L_V^{(d)}(z)=L_{V_0}^{(k)}(z')\otimes I
             +I\otimes L_0^{(d-k)}(z'').
\]
If the two $k$-dimensional matrices are isospectral for every $z'$, their
Kronecker sums with the same second factor are isospectral, proving one
direction.  Conversely, equality of the two $d$-dimensional characteristic
polynomials implies, for every integer $m\geq1$, equality of power traces
and hence of trace exponentials near $t=0$.  Using
\[
 \tr e^{t(X\otimes I+I\otimes B)}
   =(\tr e^{tX})(\tr e^{tB}),
\]
the common factor $\tr e^{tB}$ can be cancelled near zero because its value
at $t=0$ is $3^{d-k}$.  Thus the trace exponentials, and therefore the
characteristic polynomials, of $L_{V_0}^{(k)}(z')$ and $L_0^{(k)}(z')$
agree for every $z'$.
\end{proof}

\begin{corollary}[Sparse cylinder obstruction]\label{cor:cylinder}
If $d\geq3$ and a nonzero $(3\Z)^d$-periodic Floquet-isospectral potential
is, up to a coordinate permutation, a cylinder lift from three coordinate
factors, then its support in a fundamental cell satisfies
\[
 |\supp V|\geq7\cdot3^{d-3}.
\]
A cylinder depending on at most two coordinates must be zero.
\end{corollary}

\begin{proof}
For $d=3$ the first assertion is Corollary~\ref{cor:floquet}.  For $d>3$,
Proposition~\ref{prop:cylinder} and that corollary give
$|\supp V_0|\geq7$, and the lift repeats each supported value $3^{d-3}$
times.  The assertion for at most two active coordinates follows from the
small-period rigidity part of \cite[Theorem~C]{cobb2026}.
\end{proof}

\section{Reproducibility protocol}

The ancillary files are:
\begin{itemize}[leftmargin=2em]
\item \texttt{verify\_floquet\_333\_sparse.py}, for the support-four trace
table, the symbolic sixth-trace calculation, the resolvent identity and all
support-five ideals;
\item \texttt{verify\_floquet\_333\_support6.py}, for the orbit reduction,
all support-six ideals, the JSONL transcript and deterministic comparison
with a stored reference run;
\item \texttt{audit\_h333\_support6\_completeness.py}, for the independent
union-find and Burnside audit;
\item \texttt{floquet\_support6\_certificate.jsonl}, the complete hashed
support-six computation transcript;
\item \texttt{requirements.txt}, pinning the symbolic-algebra dependency;
\item the ancillary manifest, for release-file hashes and the reference
environment.
\end{itemize}
The three principal commands are
\begin{verbatim}
python verify_floquet_333_sparse.py --quiet
python audit_h333_support6_completeness.py \
  --show-collision --cross-check-five
python verify_floquet_333_support6.py --workers 6 \
  --output /tmp/floquet_support6_check.jsonl --force \
  --reference floquet_support6_certificate.jsonl
\end{verbatim}
The reference environment was Python 3.12.13 with SymPy 1.14.0
\cite{sympy}.  One reference run took about $240$ seconds for support five
and $912$ seconds for support six; runtimes are not part of the result.
Mathematical verification consists of checking the symbolic trace and
resolvent identities, reproducing the orbit and type counts, recomputing
every reduced basis, and comparing every deterministic type record with the
stored transcript.  The comparison deliberately ignores process identifiers,
platform strings and timings.  The JSONL file is an exact computation
transcript, not a standalone formal Nullstellensatz certificate containing
multiplier identities; independent verification therefore reruns the
Gr\"obner computations.

\Needspace{8\baselineskip}
\section{Scope and open directions}

The exact status is:
\begin{table}[htbp]
\centering
\caption{Proved and open strata.}
\label{tab:status}
\begin{tabular}{>{\raggedright\arraybackslash}p{0.43\textwidth}
                >{\raggedright\arraybackslash}p{0.43\textwidth}}
\toprule
stratum & status\\
\midrule
support at most $3$, every $d\geq1$ & impossible by traces $1$--$3$\\
support $4$, every $d\geq1$ & impossible by traces through order $6$\\
support $5$ in $(3\Z)^3$ & exact computer-assisted impossibility\\
support $6$ in $(3\Z)^3$ & exact computer-assisted impossibility\\
nonzero periodic-fibre shift on $H(3,3)$ & exists, necessarily support $\geq7$\\
exact support $7$ in $H(3,3)$ & open\\
unrestricted $(3,3,3)$ existence & open\\
\bottomrule
\end{tabular}
\end{table}

The nearest finite continuation is support seven, but its orbit count and
seven-variable Rabinowitsch ideals are materially larger.  A potentially more
global approach is to combine several roots-of-unity quasimomenta.  The
present proof deliberately uses only $z=1$; a few nonzero phases may impose
incompatible low-rank resolvent identities simultaneously and could obstruct
many support sizes at once.  Neither continuation is asserted here.  Nor does
the $d=3$ computation exclude arbitrary supports five or six in dimensions
$d>3$; only the cylinder consequence in Corollary~\ref{cor:cylinder} is
claimed across dimensions.

\section*{Data and code availability}

The exact scripts and the full JSONL transcript are supplied with the source
archive accompanying this manuscript.  No proprietary solver or
floating-point computation is required.

\begin{thebibliography}{99}

\bibitem{cobb2026}
J.~Cobb, M.~Faust and A.~Kretschmer,
\emph{Inverse eigenvalue problems, Floquet isospectrality and the
Hilbert--Chow morphism},
arXiv:2601.09604 (2026).

\bibitem{cox}
D.~Cox, J.~Little and D.~O'Shea,
\emph{Ideals, Varieties, and Algorithms},
Springer, Cham, 2015.

\bibitem{faustideal2026}
M.~Faust, L.~Friedman, G.~O'Malley, R.~Ramos and A.~Sharma,
\emph{Algebraic properties of the ideal of spectral invariants for the
discrete Laplacian},
arXiv:2602.08304 (2026).

\bibitem{faust2024}
M.~Faust, W.~Liu, R.~Matos, J.~Plute, J.~Robinson, Y.~Tao, E.~Tran and
C.~Zhuang,
\emph{Floquet isospectrality of the zero potential for discrete periodic
Schr\"odinger operators},
J. Math. Phys. \textbf{65} (2024), 073501.

\bibitem{kappeler1989}
T.~Kappeler,
\emph{Isospectral potentials on a discrete lattice. III},
Trans. Amer. Math. Soc. \textbf{314} (1989), 815--824.

\bibitem{liu2023}
W.~Liu,
\emph{Floquet isospectrality for periodic graph operators},
J. Differential Equations \textbf{374} (2023), 642--653.

\bibitem{saburova2026}
N.~Saburova,
\emph{Spectral invariants for discrete Schr\"odinger operators on periodic
graphs},
arXiv:2504.03282v3 (2026).

\bibitem{shipman2026}
S.~P.~Shipman and F.~Sottile,
\emph{Algebraic aspects of periodic graph operators},
in \emph{Operator Theory}, Springer, Cham, 2026, pp.~3513--3553;
arXiv:2502.03659v2.

\bibitem{sympy}
A.~Meurer et al.,
\emph{SymPy: symbolic computing in Python},
PeerJ Computer Science \textbf{3} (2017), e103.

\end{thebibliography}

\end{document}
